\documentclass[12pt]{article}

\usepackage[utf8]{inputenc}
\usepackage[T1]{fontenc}
\usepackage{amsmath,amssymb}
\usepackage{booktabs}
\usepackage{graphicx}
\usepackage{multirow}
\usepackage{xcolor}
\usepackage{url}
\usepackage{hyperref}
\usepackage{cleveref}
\usepackage{lineno}
\usepackage[round,authoryear]{natbib}
\usepackage[a4paper,margin=1in]{geometry}

\hypersetup{colorlinks=true,linkcolor=blue,citecolor=blue,urlcolor=blue}

\newcommand{\model}{AetherArgus}
\newcommand{\challenge}{VIguard-LAN}
\newcommand{\commitid}{\texttt{825fb95f\allowbreak b87a08d1\allowbreak 4e15e590\allowbreak 78c7cf88\allowbreak beb9c5e8}}
\newcommand{\labelhash}{\texttt{13fd71f5\allowbreak 507ec9bf\allowbreak 6db74c38\allowbreak fafb073c\allowbreak 76e9766c\allowbreak 0e20e874\allowbreak 452d3c62\allowbreak 302ad8ad}}
\newcommand{\figroot}{../../../ml/reports/lan-challenge-evaluation/figures/}
\newcommand{\evidencefigure}[4]{%
  \begin{figure}[t]
    \centering
    \IfFileExists{\figroot#1}{\includegraphics[width=#2\linewidth]{\figroot#1}}{\fbox{\parbox[c][0.16\textheight][c]{0.92\linewidth}{\centering Evidence figure generated in the repository:\\ \texttt{#1}}}}
    \caption{#3}
    \label{#4}
  \end{figure}}

\title{Can Compact Language Models Reliably Support HTTP Security Decisions?\\A Leakage-Controlled Empirical Evaluation}
\author{Anonymous manuscript}
\date{}

\begin{document}
\maketitle

\begin{abstract}
Compact language models are increasingly considered for local cybersecurity analysis, yet their apparent performance can be distorted by label leakage, internal test sets, binary task simplification, and inconsistent inference pipelines. This paper reports a leakage-controlled empirical evaluation of a 4B-parameter QLoRA-adapted model for request-side HTTP security decision support. The study uses a frozen 413-sample authorized LAN challenge set excluded from training, prompt selection, threshold selection, and model selection. Eight systems are compared on identical model-visible inputs, including repository rules, a character n-gram classifier, untuned zero-shot and few-shot language-model baselines, LoRA and merged artifacts, and two GGUF serving paths. The adapted model reaches 0.4658 verdict macro-F1 versus 0.4175 for zero-shot Qwen3-4B; the paired difference is 0.0482 with a 95\% bootstrap confidence interval of 0.0113--0.0835. Attack recall reaches 0.9072, but the model falsely labels 43 of 209 safe requests as attacks and identifies none of 107 insufficient-context cases as unknown. Message serialization changes 11 of 413 verdicts, while byte-identical prompts remove the observed cross-runtime verdict differences. The results show that compact models can provide attack-sensitive triage signals, but high false-positive cost and complete abstention failure make autonomous blocking unjustified. The study contributes a reproducible characterization of reliability limits, rather than a new learning architecture.
\end{abstract}

\noindent\textbf{Keywords:} HTTP security; cyberattack detection; compact language models; empirical evaluation; false positives; abstention; reproducibility

\linenumbers

\section{Introduction}
\label{sec:introduction}

HTTP security decisions increasingly involve evidence that is difficult to reduce to a single signature. A request may contain a recognizable injection string, a potentially sensitive object identifier, ambiguous authorization context, or natural-language content intended for an AI-enabled application. Signature engines remain fast and auditable, but their coverage is bounded by encoded rules. Learned classifiers can generalize beyond exact signatures, yet their reliability depends on the data distribution, the available context, and the cost assigned to false alarms. Instruction-following language models add another option because they can jointly interpret paths, headers, bodies, and contextual fields while returning a structured security judgment. The same instruction-following behavior also creates risk: attacker-controlled content may attempt to redirect the classifier, and a fluent answer may conceal an unjustified decision~\citep{perez2022promptinject,greshake2023indirect,schulhoff2023hackaprompt}.

Recent cybersecurity research has begun to evaluate local language models as traffic-analysis aids rather than assuming that general language competence transfers directly to attack detection. For example, a recent study in the \emph{Journal of Cyber Security Technology} found that local models could recognize some language-based attack indicators but remained unreliable across packet formats, attack types, and benign traffic~\citep{paukkeri2026local}. That result motivates a more specific question at the HTTP decision boundary: \emph{when a compact model is adapted to a structured request-classification task, which improvements are real, and which deployment-critical failures remain?}

This paper addresses that question through an empirical study of a 4B-parameter model adapted with QLoRA~\citep{hu2021lora,dettmers2023qlora}. The objective is not to introduce a new learning algorithm or model architecture. Instead, the study characterizes the reliability limits of compact language models used for request-side HTTP security decision support. This distinction matters because a statistically significant improvement over an untuned base model does not by itself establish operational usefulness. A model may gain attack recall while producing too many false positives, may satisfy a JSON schema while making the wrong security decision, or may force insufficient-context cases into a binary verdict instead of abstaining.

The evaluation therefore treats leakage control and failure decomposition as part of the cybersecurity evidence. The 413-sample \challenge{} challenge set is external to all training records and few-shot examples. It is excluded from prompt selection, threshold selection, hyperparameter tuning, and model selection. Prediction and scoring are separated: inference reads only model-visible request evidence, writes raw outputs, and terminates before the scoring process loads frozen labels. The study evaluates \texttt{safe}, \texttt{attack}, and \texttt{unknown} as distinct classes and reports safe false-positive rate, attack recall, abstention behavior, schema validity, latency, memory, artifact size, and paired statistical tests. Multiple model artifacts and message-serialization paths are tested because security behavior can change even when nominal weights and generation settings appear equivalent.

The principal result is deliberately mixed. The adapted model improves verdict macro-F1 from 0.4175 for zero-shot Qwen3-4B to 0.4658, with a paired difference of 0.0482 and a 95\% bootstrap confidence interval of 0.0113--0.0835. Attack recall reaches 0.9072. However, the model falsely labels 43 of 209 safe requests as attacks and predicts no unknown verdict for any of 107 insufficient-context cases. A change from separate system/user messages to one concatenated message alters 11 of 413 verdicts. These findings support use as an attack-sensitive shadow signal or analyst triage aid, but they contradict an autonomous-blocking interpretation.

The paper makes four empirical contributions:
\begin{enumerate}
    \item \textbf{A leakage-controlled external evaluation.} The study binds a frozen request-only challenge set, separated inference and scoring processes, content hashes, raw outputs, and paired statistical analysis into one auditable evaluation record.
    \item \textbf{A same-input comparison across security decision systems.} Eight systems are evaluated on identical model-visible requests, including repository rules, a character n-gram classifier, zero-shot and few-shot base-model prompts, the QLoRA adapter, merged artifacts, and GGUF serving paths.
    \item \textbf{A decomposition of cybersecurity failure modes.} The analysis separates attack detection, safe-request false positives, unknown-class abstention, structured-output validity, prompt-injection behavior, quantization effects, and message-serialization sensitivity.
    \item \textbf{Operationally bounded conclusions.} The results identify where compact models may assist triage and where they remain unsafe, particularly when false-positive cost and insufficient context are ignored.
\end{enumerate}

The remainder of the paper reviews related cybersecurity and language-model studies, defines the task and model, documents the frozen data and evaluation protocol, presents comparative and ablation results, and discusses the implications for security decision support.

\section{Related Work}
\label{sec:related}

\subsection{Rule-Based and Learned Intrusion Detection}
Signature-oriented network security systems emphasize transparency and deterministic protocol analysis. Snort operationalized lightweight signature matching~\citep{roesch1999snort}, while Bro/Zeek demonstrated stateful network analysis and policy-driven event processing~\citep{paxson1999bro}. These systems remain important because their behavior is auditable and known-pattern false positives can be controlled. Their limitation is coverage: semantic variants, authorization-sensitive requests, and business-logic abuse may not be reducible to a fixed payload pattern.

Machine-learning systems broaden the decision boundary beyond explicit signatures. DeepLog models execution logs as sequences for anomaly detection~\citep{du2017deeplog}; LUCID presents a lightweight deep-learning approach to DDoS classification~\citep{doriguzzi2020lucid}; and contextual encoders have been applied to network and intrusion-detection features~\citep{yang2025bertids}. Surveys of explainable intrusion detection note recurring trade-offs between broader pattern recognition, interpretability, and distributional stability~\citep{neupane2022xids}. The present study focuses on request-side HTTP evidence and a typed security decision rather than a single anomaly score.

\subsection{Local Language Models for Cybersecurity Analysis}
Language models are increasingly studied as aids for security triage, log analysis, vulnerability interpretation, and network-traffic reasoning. Their appeal is that heterogeneous evidence can be represented as text and mapped to structured explanations or decisions. Their weakness is that fluent output can overstate certainty, and models may respond differently to minor changes in prompt structure or data representation.

The closest recent empirical precedent is the evaluation of local LLMs for network traffic packet analysis by \citet{paukkeri2026local}. That study examined Llama 3.1 in depth and compared five local models on IoT packet data. It found that models recognized some language-based malicious indicators, especially SQL-related content, but were unreliable across attack types, benign traffic, hexadecimal input, and encrypted payloads. The present work shares the reliability-oriented perspective but differs in five ways. First, it studies request-level HTTP security decisions rather than individual packet interpretation. Second, it evaluates a task-adapted compact model against rules, conventional machine learning, untuned prompts, merged artifacts, and multiple serving paths. Third, it uses a frozen external challenge set that is isolated from training and tuning. Fourth, it treats insufficient context as a scored \texttt{unknown} class. Fifth, it applies paired confidence intervals and significance tests to same-sample predictions.

The contribution is therefore not another claim that an LLM can recognize obvious attack strings. It is a controlled estimate of how much task adaptation changes decisions, which failure modes remain, and whether those changes support a defensible security role.

\subsection{Parameter-Efficient Adaptation and Structured Decisions}
Transformers provide the underlying sequence-modeling architecture~\citep{vaswani2017attention}. The evaluated base model is Qwen3-4B-Instruct-2507~\citep{qwen3,qwen2507}. LoRA freezes pretrained weights and learns low-rank updates~\citep{hu2021lora}; QLoRA combines low-rank adapters with a frozen 4-bit quantized base model~\citep{dettmers2023qlora}. These techniques make local specialization feasible, but they do not guarantee calibrated uncertainty, safe false-positive behavior, or domain transfer.

This paper uses QLoRA as an established adaptation mechanism, not as a methodological novelty. The research question is empirical: whether adaptation improves attack-oriented HTTP classification under a sealed protocol, and whether the resulting error profile remains compatible with security decision support.

\subsection{Prompt Injection and Untrusted Request Content}
Prompt injection exploits the inability of instruction-following models to reliably separate trusted control text from untrusted content. PromptInject demonstrated goal hijacking and prompt leakage~\citep{perez2022promptinject}; indirect prompt injection extended the threat to retrieved or tool-supplied content~\citep{greshake2023indirect}; and HackAPrompt documented diverse bypass strategies~\citep{schulhoff2023hackaprompt}. In an HTTP classifier, the body, path, query string, and some headers are attacker-controlled data. The model must analyze those fields without obeying embedded instructions. The study therefore preserves raw outputs and evaluates message structure, schema compliance, and prompt-injection cases separately from verdict accuracy.

\subsection{Evaluation, Abstention, and Reproducibility}
Security evaluation is vulnerable to leakage, prevalence distortion, and selective metric reporting. We use stratified bootstrap confidence intervals~\citep{efron1993bootstrap}, paired McNemar tests~\citep{mcnemar1947note}, and Holm--Bonferroni correction~\citep{holm1979sequentially}. In addition to accuracy and F1, we report safe false-positive rate and unknown-class recall. The unknown class represents cases in which request-visible evidence is insufficient for a justified safe/attack decision; operationally, it functions as an abstention target.

A model can fail by choosing the wrong verdict, by following attacker-controlled instructions, by producing an invalid schema, or by expressing unjustified certainty. Reporting only aggregate accuracy collapses these failure channels. The present work scores all samples and reports each channel explicitly, while also testing whether artifact format and message serialization alter decisions.

\section{Task Formulation and Model}
\label{sec:method}

\subsection{Request-Side Classification Task}
Let an HTTP request be represented as
\begin{equation}
    x = (m, p, h, b, c),
\end{equation}
where $m$ is the method, $p$ the path and query, $h$ the model-visible headers, $b$ the request body, and $c$ a bounded source-context description available at decision time. Response status, response body, hidden object ownership, collector labels, and adjudication metadata are excluded from inference.

The primary prediction is
\begin{equation}
    \hat{y} \in \{\texttt{safe},\texttt{attack},\texttt{unknown}\}.
\end{equation}
The \texttt{unknown} label is not a miscellaneous attack category. It denotes insufficient model-visible evidence for a justified safe/attack decision. This distinction is operationally important because authorization, IDOR/BOLA, business-logic, and workflow attacks can be impossible to resolve from request syntax alone.

Each generative model must return a JSON object containing:
\begin{itemize}
    \item \texttt{verdict}: \texttt{safe}, \texttt{attack}, or \texttt{unknown};
    \item \texttt{attackFamily}: a fixed enumeration including injection, authorization, traversal, SSRF, prompt injection, and related families;
    \item \texttt{recommendedAction}: a bounded action recommendation;
    \item \texttt{contextSufficient}: whether request-side evidence is adequate;
    \item \texttt{confidence}: an ordinal confidence value;
    \item supporting fields required by the frozen schema.
\end{itemize}
Raw model text is preserved before parsing. Invalid JSON, missing fields, and illegal enum values are counted separately. A parse failure is mapped to \texttt{unknown} for scoring rather than silently repaired to \texttt{safe}.

\subsection{Base Model and Adaptation}
\model{} v0.2 is derived from Qwen3-4B-Instruct-2507. The model is adapted using QLoRA: the base weights are loaded in 4-bit NF4 form and kept frozen while rank-16 low-rank adapters are trained with scaling factor 32 and dropout 0.05. Training uses BF16 computation where supported. The locked training seed is 20260728.

The training run used a single NVIDIA RTX 5070 Ti Laptop GPU with 11.94~GiB available memory. One epoch required 12,695 seconds (approximately 3.53 hours) and reached a peak GPU allocation of 6.77~GiB. The run completed 895 optimization steps with reported training loss 0.07975 and validation loss 0.00963. Only the adapter parameters are updated; the resulting adapter occupies approximately 132~MB, while the combined base-plus-adapter directory measures approximately 7.62~GiB.

The system prompt defines the evidence boundary, output schema, enum constraints, and instruction hierarchy. It explicitly treats all request content as untrusted data. For the frozen challenge evaluation, all generative systems use the same prompt bytes, a 2,048-token context cap, a 320-token output cap, deterministic decoding, and seed 0 where the runtime supports it.

\subsection{Training Objective and Data Mixture}
Training examples serialize request-side evidence and a target JSON decision. The corpus combines deterministic synthetic generation and public-source-derived examples, including normal requests, attack families, hard negatives, unknown-context cases, and prompt-injection attempts. The objective is standard causal language modeling over the complete response sequence. Because schema tokens and verdict tokens are learned jointly, task adaptation can improve both semantic classification and serialization compliance. The evaluation therefore reports these components separately.

The v0.2 release contains 14,314 training examples, 1,752 validation examples, 1,236 closed-test examples, and 1,750 sealed independent examples. Split construction is deterministic and audited for exact and near-duplicate leakage. The sealed independent split was created before later development generation. None of these internal test sets are used to select the final external-challenge prompt or threshold.

\subsection{Inference Artifacts}
We evaluate four representations of the adapted model:
\begin{enumerate}
    \item \textbf{HF-LoRA}: Qwen3-4B plus the learned adapter through Transformers/PEFT;
    \item \textbf{Merged BF16}: adapter weights merged into a BF16 Transformers model;
    \item \textbf{Ollama f16}: an f16 GGUF artifact served through Ollama;
    \item \textbf{GGUF Q4\_K\_M}: a quantized GGUF artifact for reduced local storage and faster tested inference.
\end{enumerate}
Measured artifact sizes are 8.18~GB for the base-plus-adapter directory, 8.04~GB for merged BF16, 8.05~GB for f16 GGUF, and 2.50~GB for Q4\_K\_M (decimal units). These artifacts share a model lineage but are not assumed to be numerically identical.

\subsection{Safety Boundary}
\model{} does not directly execute a block, delete, credential, or hardware action. A model output is an analytical signal consumed by a separate policy layer. The intended deployment mode for the evaluated version is shadow analysis or second-stage review. This boundary is essential because a high-recall classifier with a 20.57\% safe false-positive rate is unsuitable as an autonomous enforcement authority.

\section{Datasets, Labels, and Leakage Controls}
\label{sec:data}

\subsection{Internal Development Data}
The internal v0.2 corpus supports training, validation, and development-time evaluation. Its examples are assembled from deterministic templates and public-source-derived security cases. The data mixture deliberately includes hard negatives---requests that contain suspicious tokens but should remain safe---and an explicit unknown class for cases in which model-visible context is insufficient. Prompt-injection examples target both the verdict and the recommended-action fields.

Internal evaluation is useful for checking training convergence and output alignment, but it is not the primary evidence in this paper. On the sealed 1,750-example independent split, the LoRA model achieved 0.9657 verdict accuracy, 1.000 malicious recall, 0.0286 normal FPR, and 1.000 unknown accuracy. It nevertheless failed three locked release gates: JSON syntax success was 0.9714, schema validity was 0.9714, and prompt-injection obedience was 0.9197. Eleven of 137 prompt-injection cases hijacked the verdict or action recommendation. These results motivated an external challenge evaluation rather than a production-readiness claim.

\subsection{Frozen \challenge{} Challenge Set}
The primary test set contains 413 authorized LAN request-side samples. Its model-input file has SHA-256
\begin{center}
{\small\texttt{c77f4949\allowbreak cbe1860f\allowbreak 97219a83\allowbreak f0ee16de\allowbreak 40740321\allowbreak a4d2f628\allowbreak d56f1506\allowbreak 9fe6b8c3}}.
\end{center}
The frozen reference-label distribution is 209 safe, 97 attack, and 107 unknown. The label file has SHA-256 \labelhash. The set is external to the v0.2 train, validation, closed-test, and independent-test splits. It is not used for training, prompt selection, threshold selection, hyperparameter tuning, checkpoint selection, or model selection.

The challenge set contains request-side fields only. Labels and scoring metadata are stored separately. Cohorts include normal requests, attack attempts, hard negatives, prompt-injection cases, and insufficient-context examples. The dataset is an authorized challenge capture, not a random sample of natural production traffic; prevalence-sensitive metrics must therefore not be interpreted as production base rates.

\subsection{AI-Assisted Annotation and Adjudication}
Two independent AI annotators produced initial labels. Agreement was high but imperfect, with observed agreement 0.9274 and Cohen's $\kappa=0.8651$. A third AI model adjudicated disagreements. A separate model-visible-context audit corrected 68 provenance issues in which labels depended on information that the classifier could not observe. The resulting frozen reference set contains 160 direct-consensus outcomes and 253 model-adjudicated outcomes. The reference file was frozen before the AetherArgus challenge predictions were generated.

The AI-assisted stage is not described as human ground truth. It is a reproducible label-construction process with known dependence on model judgment. To reduce this threat, the full frozen reference set was subsequently subjected to a confidential blind human review.

\subsection{Confidential Blind Human Review}
Human review occurred after the frozen model evaluation. Reviewers operated outside the public repository in a restricted environment and made their decisions before seeing annotator-A or annotator-B labels, the third-model adjudication, AetherArgus predictions, or evaluation metrics. After blind decisions were complete, the human outcomes were compared with the frozen references.

The review covered all 413 samples, found zero disagreements, and required zero label changes. The before- and after-review label hashes are therefore identical. No prediction rerun or rescoring was required. The repository records \texttt{humanReviewed=true}, full coverage, zero disagreements, and a SHA-bound operator attestation.

Reviewer identities, reviewer count, worksheets, and raw materials are withheld because they contain confidential benchmark and operational information. This limits external auditability. We consequently use the term \emph{confidentially human-reviewed reference labels}, not publicly reproducible human ground truth.

\subsection{Frozen Protocol and Process Isolation}
Protocol v1 was committed before the rule baseline. Protocol v2 was committed before the remaining seven primary model runs and fixes model identities, data hashes, metrics, statistical tests, system-prompt hash, few-shot example hash, model-asset hashes, and deterministic generation parameters.

Inference and scoring are process-separated:
\begin{equation}
    \text{model inputs} \rightarrow \text{predictions.jsonl}
\end{equation}
\begin{equation}
    \text{predictions.jsonl} + \text{labels} \rightarrow \text{metrics}.
\end{equation}
The inference programs do not open the label file or adjudication artifacts. Predictions are append-only, flushed per sample, and keyed by unique sample identifiers. The strict validator checks 413-of-413 coverage, duplicate and omission absence, common input hash, prediction-label alignment, train/test separation, few-shot exclusion, metric recomputation, figure/table consistency, privacy scans, and claim-to-evidence traceability. The frozen evidence package passed all 27 strict checks.

\section{Experimental Setup}
\label{sec:setup}

\subsection{Research Questions}
The evaluation addresses the following questions:
\begin{description}
    \item[RQ1] Does task-adapted \model{} outperform the untuned 4B base model and conventional baselines on the frozen external challenge set?
    \item[RQ2] Are gains attributable only to structured-output alignment, or do verdict-level security decisions improve?
    \item[RQ3] How does the model trade attack recall against false positives and abstention?
    \item[RQ4] What quality and storage trade-offs arise from merged and quantized artifacts?
    \item[RQ5] Do message structure, prompt serialization, or runtime path change verdicts?
\end{description}

\subsection{Compared Systems}
All eight primary systems receive the same 413 request-side samples.

\textbf{Rules baseline.} A read-only offline adapter directly invokes the repository's frozen detection engine. It does not access response evidence or labels and is not extended after challenge outcomes are observed.

\textbf{Character n-gram logistic regression.} This baseline uses character 3--5 grams over serialized request-side evidence, a maximum of 250,000 features, minimum document frequency 2, balanced class weights, and deterministic random seed 20260728. Training uses only v0.2 train and validation data.

\textbf{Qwen3-4B zero-shot.} The base instruction model receives the frozen task description and output schema without demonstrations.

\textbf{Qwen3-4B fixed few-shot.} The base model receives a precommitted set of demonstrations selected from train/validation data. Challenge examples are excluded, and the example file is hash-locked before evaluation.

\textbf{\model{} artifacts.} We evaluate HF-LoRA, merged BF16, Ollama f16 GGUF, and GGUF Q4\_K\_M variants. The primary analysis uses HF-LoRA; the others measure merge, runtime, and quantization effects.

\subsection{Metrics}
The primary metric is three-class Macro-F1, which prevents the large safe class from dominating the reported score. We also report accuracy, weighted F1, per-class precision/recall/F1, and a $3\times3$ confusion matrix. Attack detection is additionally evaluated as attack versus safe-or-unknown and as attack versus safe only.

Safe false-positive rate is
\begin{equation}
 \mathrm{FPR}_{safe}=\frac{\#(y=\texttt{safe},\hat y=\texttt{attack})}{\#(y=\texttt{safe})}.
\end{equation}
Unknown performance includes precision, recall, accuracy on reference-unknown samples, and unknown-to-safe/attack transition rates. Attack-family Macro-F1 is computed only over attack families present in the reference attack subset; a predicted \texttt{none} counts as an error but is not added as a zero-support reference class.

Structured-output quality includes JSON syntax success, complete-field rate, enum validity, and schema validity. Performance evidence includes per-sample latency, throughput, measured artifact size, and available memory observations. Cross-runtime latency is descriptive because timing boundaries and serving stacks differ.

\subsection{Statistical Analysis}
Point estimates are accompanied by 10,000-iteration stratified bootstrap 95\% confidence intervals with seed 20260731. Pairwise verdict correctness is compared using McNemar's test, using the exact form when discordant counts are small. Paired stratified bootstrap estimates differences in Macro-F1, attack F1, safe FPR, and unknown accuracy. Holm--Bonferroni correction controls family-wise error across planned model comparisons. For binomial proportions, Wilson and Clopper--Pearson intervals are retained in the machine-readable evidence, particularly for zero-event outcomes where bootstrap intervals can degenerate.

\subsection{Hardware and Software}
Primary generative-model runs were executed on Windows 11 with an NVIDIA GeForce RTX 5070 Ti Laptop GPU. The recorded environment includes Python 3.12.10, a CUDA 13.0 development stack, PyTorch 2.12.0 development build, Transformers 5.2.0, PEFT 0.20.0, and Ollama 0.32.5. Exact model revisions, prompt hashes, runtime settings, and artifact hashes are stored in the evaluation manifest. This unusual forward-version environment is reported for reproduction rather than treated as a performance advantage.

\subsection{Inference-Path Diagnostics}
Four pre-registered diagnostics isolate message and serialization effects:
\begin{enumerate}
    \item merged BF16 with native \texttt{system+user} messages;
    \item merged BF16 with a single concatenated message;
    \item Ollama f16 with native \texttt{system+user} messages;
    \item Ollama f16 with the same HF-pre-serialized prompt bytes used by the Transformers path.
\end{enumerate}
Equivalent comparisons are repeated with shared serialized prompts. Each diagnostic covers all 413 samples. Because Transformers uses a BF16 merged artifact while Ollama uses an f16 GGUF artifact, residual differences cannot be attributed solely to runtime even when prompt bytes are identical.

\section{Results}
\label{sec:results}

\subsection{Main Comparative Results}
\Cref{tab:main-results} reports the eight-system comparison. HF-LoRA reaches 0.6150 accuracy and 0.4658 Macro-F1. Merged BF16 and Ollama f16 produce the same accuracy and nearly identical Macro-F1. Q4\_K\_M loses 0.0102 Macro-F1 relative to merged BF16 while retaining 0.9072 attack recall.

\begin{table*}[t]
\centering
\caption{Results on the frozen 413-sample \challenge{} challenge set. FPR is computed only over the 209 reference-safe samples. Unknown is accuracy on 107 reference-unknown samples.}
\label{tab:main-results}
\resizebox{\textwidth}{!}{%
\begin{tabular}{lrrrrrrr}
\toprule
System & Accuracy & Macro-F1 & Attack F1 & Attack Recall & Safe FPR & Unknown & Schema Valid \\
\midrule
Frozen rules & 0.5206 & 0.2994 & 0.1606 & 0.1134 & \textbf{0.0144} & 0.0187 & 1.0000 \\
Char n-gram LR & 0.2349 & 0.1273 & 0.3819 & \textbf{1.0000} & 1.0000 & 0.0000 & 1.0000 \\
Qwen3-4B zero-shot & 0.5230 & 0.4175 & 0.4964 & 0.7113 & 0.3206 & 0.0561 & 1.0000 \\
Qwen3-4B fixed few-shot & 0.5133 & 0.4114 & 0.4758 & 0.6082 & 0.2679 & \textbf{0.0654} & 0.9976 \\
\model{} HF-LoRA & \textbf{0.6150} & 0.4658 & 0.6692 & 0.9072 & 0.2057 & 0.0000 & 1.0000 \\
\model{} merged BF16 & \textbf{0.6150} & \textbf{0.4661} & \textbf{0.6718} & 0.9072 & 0.2057 & 0.0000 & 1.0000 \\
\model{} Ollama f16 & \textbf{0.6150} & \textbf{0.4661} & \textbf{0.6718} & 0.9072 & 0.2057 & 0.0000 & 1.0000 \\
\model{} Q4\_K\_M & 0.6005 & 0.4559 & 0.6519 & 0.9072 & 0.2249 & 0.0000 & 0.9952 \\
\bottomrule
\end{tabular}}
\end{table*}

\evidencefigure{figure-1-model-comparison.pdf}{0.98}{Model comparison generated from the machine-readable frozen metrics.}{fig:model-comparison}

The character n-gram classifier predicts every challenge sample as attack. This yields perfect attack recall but 100\% safe FPR and zero unknown accuracy. The behavior is evidence of severe domain transfer failure rather than a useful operating point. The frozen rules baseline exhibits the opposite profile: safe FPR is only 3/209 (1.44\%), but it identifies 11/97 attacks (11.34\% recall). \model{} therefore does not dominate the rules baseline in every operational dimension; it exchanges precision for recall.

\subsection{Significance of Task Adaptation}
Relative to Qwen3-4B zero-shot, HF-LoRA improves Macro-F1 by 0.0482 with paired-bootstrap 95\% CI [0.0113, 0.0835]. Attack F1 improves by 0.1728 with 95\% CI [0.1246, 0.2253]. Safe FPR decreases by 0.1148 with 95\% CI [0.0718, 0.1579] in absolute magnitude. McNemar's test on verdict correctness remains significant after Holm correction ($p=8.36\times10^{-7}$).

Against the fixed few-shot baseline, HF-LoRA improves Macro-F1 by approximately 0.0544 with 95\% CI [0.0157, 0.0912]; the Holm-corrected McNemar value is $7.22\times10^{-7}$. Fixed demonstrations therefore do not reproduce the benefit of task adaptation. In fact, few-shot attack recall falls to 0.6082, below the zero-shot value of 0.7113.

The verdict improvement is not merely a schema effect. Both Qwen zero-shot and HF-LoRA achieve 100\% schema validity on the challenge set, yet their verdict metrics differ significantly. This supports RQ2: the adapter changes security decisions, not only serialization.

\subsection{False Positives and Abstention Failure}
The central negative result is visible in the HF-LoRA confusion matrix:
\begin{center}
\begin{tabular}{lrrr}
\toprule
Reference $\backslash$ Prediction & Safe & Attack & Unknown \\
\midrule
Safe & 166 & 43 & 0 \\
Attack & 9 & 88 & 0 \\
Unknown & 72 & 35 & 0 \\
\bottomrule
\end{tabular}
\end{center}
The model detects 88 of 97 attacks but incorrectly flags 43 of 209 safe requests. More importantly, it never emits \texttt{unknown}. All 107 insufficient-context samples are forced into safe (72) or attack (35). Compared with Qwen zero-shot, unknown accuracy decreases by 0.0561 with paired-bootstrap 95\% CI [-0.1028,-0.0187]. Task adaptation therefore sharpens the attack decision boundary while degrading abstention.

\evidencefigure{figure-2-verdict-confusion-matrices.pdf}{0.98}{Verdict confusion matrices. The adapted model's high attack recall coexists with complete unknown-class collapse.}{fig:confusion}

These observations rule out a production-safe autonomous-blocking claim. A classifier that flags roughly one in five reference-safe requests and cannot identify insufficient evidence should not directly control irreversible actions.

\subsection{Attack-Family Classification}
On reference-attack samples, HF-LoRA reaches attack-family Macro-F1 0.6425. Qwen zero-shot reaches 0.6325, while fixed few-shot reaches 0.6661. The family-label result is therefore weaker than the verdict result: task adaptation improves whether a request is classified as an attack, but does not establish superiority in fine-grained family naming. Family metrics are also based on small per-family supports and are treated as secondary evidence.

\subsection{Quantization and Artifact Size}
\Cref{tab:size} summarizes measured artifact sizes and descriptive latency. Q4\_K\_M reduces storage to 2.50~GB (2.33~GiB), approximately 31\% of the merged-BF16 directory size. Its Macro-F1 is 0.4559 compared with 0.4661 for merged BF16. The paired difference of 0.0102 has 95\% CI [-0.0043,0.0249], and the Holm-corrected McNemar comparison is not significant ($p=0.539$). We therefore conclude that Q4 preserves most verdict-level quality on this set, not that it is lossless.

\begin{table}[t]
\centering
\caption{Measured artifact size and descriptive latency. Runtime stacks and timing boundaries differ; latency values are not a controlled speed ranking.}
\label{tab:size}
\resizebox{\linewidth}{!}{%
\begin{tabular}{lrrr}
\toprule
Artifact & Size (GB) & P50 (ms) & P95 (ms) \\
\midrule
Qwen3-4B zero-shot & 8.04 & 2187 & 2620 \\
Qwen3-4B few-shot & 8.04 & 4977 & 259883 \\
HF base + LoRA & 8.18 & 2730 & 22654 \\
Merged BF16 & 8.04 & 12103 & 15056 \\
Ollama f16 GGUF & 8.05 & 1679 & 1935 \\
Q4\_K\_M GGUF & 2.50 & 446 & 1535 \\
\bottomrule
\end{tabular}}
\end{table}

\evidencefigure{figure-5-quality-latency-size-tradeoff.pdf}{0.98}{Quality, latency, and measured artifact-size trade-offs. Cross-stack latency is shown descriptively.}{fig:tradeoff}

The extreme few-shot P95 value reflects long-tail generation behavior under that path. Because model loading, HTTP serving, prompt length, and measurement boundaries differ, the paper does not claim a multiplicative runtime advantage from these observations.

\subsection{Inference Path and Message Serialization}
Changing from native \texttt{system+user} messages to a single concatenated message changes 11 of 413 verdicts in the Transformers path and 11 of 413 in the Ollama path. Safe FPR rises from 0.2057 to 0.2249 for Transformers and to approximately 0.2297 for Ollama. Message structure is therefore an isolated source of security-decision variance.

When the exact same HF-pre-serialized prompt bytes are supplied to the merged-BF16 and Ollama-f16 paths, verdict agreement is 413/413 and both produce safe FPR 43/209. Native \texttt{system+user} paths also agree on all 413 verdicts in this experiment. This is evidence that prompt-byte control can eliminate the observed verdict discrepancy for this dataset; it is not proof that the runtimes or BF16/f16 artifacts are numerically equivalent on all inputs.

\evidencefigure{figure-6-inference-path-ablation.pdf}{0.98}{Pre-registered inference-path diagnostics. Message structure changes decisions, while shared serialized prompts produce identical tested verdicts.}{fig:ablation}

\subsection{Internal-to-External Generalization Gap}
The sealed internal independent split and external challenge set produce sharply different outcomes. Internally, verdict accuracy is 0.9657, normal FPR 0.0286, and unknown accuracy 1.000. On \challenge{}, accuracy falls to 0.6150, safe FPR rises to 0.2057, and unknown accuracy falls to zero. The datasets differ in composition and cannot be compared as identical populations, but the direction and magnitude expose substantial domain shift. Internal gate performance alone would have overstated deployment readiness.

The gap also reveals a likely training imbalance: the model learns a strong mapping from suspicious request patterns to attack, yet does not preserve the intended relationship between \texttt{contextSufficient=false} and \texttt{verdict=unknown} under external conditions. This finding motivates explicit abstention calibration and human-reviewed unknown examples in future training.

\section{Discussion and Threats to Validity}
\label{sec:discussion}

\subsection{Statistical Improvement Is Not Operational Readiness}
The adapted model improves verdict macro-F1 and attack recall relative to the untuned base model, and the paired confidence interval for macro-F1 excludes zero. This establishes that task adaptation changes same-sample decisions in a favorable attack-oriented direction. It does not establish that the resulting system is safe to deploy as a blocking authority.

The distinction is visible in the class-specific results. A safe false-positive rate of 20.57\% would impose substantial operational cost in a blocking or high-volume alerting system. Unknown recall of zero is more serious than an ordinary class imbalance: every insufficient-context case is forced into a substantive safe or attack judgment. The model therefore becomes more attack-sensitive while losing the ability to express that request-visible evidence is inadequate. For security decision support, this is a reliability failure rather than a cosmetic metric weakness.

This interpretation is consistent with recent local-LLM traffic analysis, which found useful recognition of some language-based attack indicators but insufficient reliability across benign traffic, representations, and attack types~\citep{paukkeri2026local}. The present results extend that concern to an adapted request-level classifier: domain tuning can improve recall without solving false-positive cost or abstention.

\subsection{What the Adapter Appears to Learn}
The QLoRA adapter appears to strengthen associations between security vocabulary, suspicious request compositions, and attack-oriented verdicts. Attack recall rises markedly relative to zero-shot inference, and attack F1 improves. The training objective, however, rewards production of a concrete target sequence and valid schema; it does not independently reward calibrated refusal when context is missing. A plausible consequence is that the model learns that a complete response should contain a decisive verdict even when the reference class is unknown.

This explanation remains an inference from behavior rather than a mechanistic claim. The study does not include activation analysis or causal representation probing. A stronger follow-up should train an explicit reject option, use contrastive pairs with sufficient and insufficient context, evaluate calibration, and separate verdict generation from uncertainty estimation.

\subsection{Rules and Language Models Are Complementary}
The frozen rules baseline and the adapted language model exhibit complementary error profiles. Rules produce low safe FPR but miss most attacks in the challenge set. The adapted model detects far more attacks but flags many safe requests. This supports a layered role rather than replacement:
\begin{enumerate}
    \item high-specificity rules handle known patterns with auditable behavior;
    \item the compact model analyzes unresolved requests in shadow mode or as a second-stage triage signal;
    \item disagreements and insufficient-context cases are escalated to policy logic or human review;
    \item irreversible actions require an independent authorization mechanism.
\end{enumerate}

The empirical contribution is the identification of these operating points under the same request evidence. It would be misleading to summarize the result only as a model improvement, because the safety cost is concentrated in the safe and unknown classes.

\subsection{Message Serialization Is Security-Relevant Configuration}
Combining trusted policy and untrusted request content into one message changes 11 of 413 verdicts and increases safe false positives. When the serialized prompt bytes are held identical, the tested Transformers and Ollama paths agree on every verdict in the frozen set. These observations indicate that message boundaries, chat templates, and serialization are part of the security-critical model configuration.

Weights, temperature, and nominal prompt text are therefore insufficient for reproducibility. A deployable evaluation record should version the exact message roles, template revision, serialized bytes or their digest, tokenizer, model artifact, and runtime. The 413-sample result does not establish universal cross-runtime equivalence, but it demonstrates that uncontrolled prompt construction can create variation that might otherwise be misattributed to quantization or serving software.

\subsection{Reference Labels and Human Review}
The frozen references were initially produced through dual-model annotation and independent model adjudication, followed by a model-visible-context audit. After evaluation, a confidential blind human review covered all 413 samples and produced no label changes. The unchanged label file means that the original predictions and paired statistics remain applicable, while the review reduces concern that obvious labeling errors explain the principal results.

The review is not equivalent to a publicly reproducible human-annotation corpus. Reviewer identities, worksheets, and per-sample rationales are withheld, so external researchers cannot independently inspect the review process. The paper therefore avoids describing the labels as objective ground truth. The temporal order is also disclosed: human review occurred after the model runs and did not influence model selection or tuning.

\subsection{Threats to Internal Validity}
The research team developed the training corpus, prompt contract, and adapted model. Unobserved design choices may favor the adapted system. The challenge set was excluded from training, prompt selection, threshold selection, and model selection, and all principal systems receive the same model-visible input, which limits but does not eliminate this risk.

The model was trained with one seed and one epoch. The reported confidence intervals quantify uncertainty over challenge samples, not training-run variance. Repeating training with multiple seeds is necessary before claiming optimization stability.

The character n-gram baseline collapses toward the attack class and is not a strong conventional competitor. The rules engine is operationally informative but repository-specific. Future work should include independently maintained WAF or IDS baselines and stronger calibrated text classifiers.

\subsection{Threats to Construct Validity}
A three-class verdict compresses a complex security decision. Authorization and business-logic cases may depend on authenticated identity, object ownership, workflow history, tenant boundaries, or response evidence that is absent from the request. The unknown class represents this limitation, but the evaluated model fails to use it.

Schema validity measures interface compliance rather than security correctness. Conversely, mapping parse failures to unknown can reward non-response. The study reports raw parse and schema outcomes so verdict metrics are not interpreted in isolation.

Attack-family labels are secondary and may simplify multi-causal requests. A sample can combine prompt injection, authorization abuse, and protocol manipulation. Family-level macro-F1 is therefore not treated as the primary endpoint.

\subsection{Threats to External Validity}
The challenge set contains 413 authorized LAN samples curated for security evaluation. It is not prevalence-matched production traffic, so its accuracy and false-positive rates cannot be interpreted as universal deployment rates. Family-specific estimates are particularly uncertain for rare classes.

The study evaluates one compact base-model family and one adaptation run. Results may not transfer to other tokenizers, model families, instruction-tuning recipes, or languages. The strongest claim is local: under the locked protocol, the adapted model improves attack-oriented classification over its specific prompting baselines while remaining unsuitable for autonomous blocking.

Runtime measurements were collected on a single consumer-laptop GPU and across serving stacks with different timing boundaries. Artifact sizes are measured exactly, but latency should be remeasured on the intended deployment platform under a synchronized benchmark design.

\subsection{Research Priorities}
The observed failures define a focused next stage:
\begin{itemize}
    \item train and calibrate an explicit abstention mechanism using human-reviewed pairs with insufficient context;
    \item optimize false-positive cost without sacrificing attack recall, including cost-sensitive objectives or a rule--model cascade;
    \item repeat adaptation across multiple random seeds;
    \item evaluate calibration, selective risk, and risk--coverage trade-offs;
    \item expand testing with independently sourced, production-like authorized HTTP traffic;
    \item compare with independently maintained security engines and stronger non-LLM classifiers;
    \item test adaptive prompt-injection and evasion attacks against the complete serialization pipeline.
\end{itemize}

\section{Conclusion}
\label{sec:conclusion}

This paper evaluated a compact QLoRA-adapted language model for request-side HTTP security decision support under a frozen, leakage-controlled external protocol. The study does not propose a new learning architecture. Its contribution is an empirical characterization of what task adaptation improves, which reliability failures remain, and how inference-path choices affect security judgments.

On 413 identical challenge requests, the adapted model improves verdict macro-F1 from 0.4175 for zero-shot Qwen3-4B to 0.4658. The paired difference is 0.0482, with a 95\% bootstrap confidence interval of 0.0113--0.0835, and attack recall reaches 0.9072. These results show that compact local models can acquire useful attack-oriented decision behavior.

The same evidence rules out a deployment-ready interpretation. The model classifies 43 of 209 safe requests as attacks and does not predict unknown for any of 107 insufficient-context cases. A change in message structure alters 11 of 413 verdicts, showing that prompt serialization is part of the security-critical configuration. The rules baseline remains substantially more precise on safe requests, although it misses most attacks.

Accordingly, the evaluated model is best treated as a shadow analyzer, second-stage triage signal, or source of analyst prioritization. It should not be granted autonomous blocking authority. The next technical priorities are explicit abstention, lower false-positive cost, calibrated selective prediction, multi-seed validation, stronger independent baselines, and broader authorized external testing. More generally, cybersecurity evaluations of language models should report class-specific failure costs and inference-path controls rather than equating aggregate improvement with operational reliability.


\section*{Data availability}
A de-identified reviewer package contains the model-visible challenge inputs, frozen reference labels, per-sample predictions, protocol files, aggregate metrics, statistical outputs, and evaluation scripts. It is intended to accompany the submission as supplementary material and to be deposited in a persistent repository subject to the removal of operationally sensitive material. Raw LAN capture material and confidential human-review worksheets are not publicly shared because they contain security-sensitive information. The paper evidence baseline is merge commit \commitid, and the frozen reference-label file is identified by SHA-256 \labelhash.

\section*{Ethics and responsible use}
All challenge traffic was generated and collected in an authorized LAN environment. The study does not claim that \model{} is suitable for autonomous blocking. The evaluated system is limited to shadow-mode analysis or second-stage triage; high-impact actions require independent policy, authorization, or human review.

\section*{Generative AI disclosure}
During manuscript preparation, generative AI tools were used for organization, language drafting, and consistency checking. The author reviewed and verified the manuscript and remains responsible for all claims, analyses, and conclusions.

\bibliographystyle{plainnat}
\bibliography{../references,sections/jcst-addendum}

\end{document}
